\section{Purified for the Vision of God}

Those who die in God's grace and friendship but remain imperfectly purified are assured of heaven. They undergo the final purification which the Church calls purgatory (CCC 1030--1032). Their friendship with God endures, and their destiny will not change into damnation. The purification completes the holiness required for the vision of God. In those already saved, the Lord removes what remains incompatible with that communion.

This teaches us to desire more than escape from condemnation. Sin can leave attachments which survive a real return to God. We may sincerely repent of pride while still craving the attention which nourished it. We may stop an injustice while retaining the habit of treating another person as inferior. We may give up a pleasure forbidden by God's law while continuing to resent the law. These attachments need healing. Grace calls us toward an undivided love in which obedience is no longer accompanied by a protected wish to be free of God.

The Church distinguishes the guilt of sin, eternal punishment, and temporal punishment. Forgiveness and the restoration of communion with God remit the eternal punishment of sin; temporal punishment can remain. Purification from sinful attachments takes place here or, for those who die in grace still needing it, after death. A conversion of fervent charity can be so complete that no punishment remains (CCC 1472--1473). The language of punishment expresses the seriousness of sin and the restoration of order. The language of healing expresses the transformation of the person. Both concern the same saving work of God.

An earthly analogy can help without measuring the mystery. Someone who has lied can be truly forgiven while the habit of lying still needs to be overcome and the truth restored to those deceived. The forgiveness is real from the beginning; its reality makes a new course possible. It does not make the falsehood good or remove the need for amendment. Purgatory concerns a far deeper purification, accomplished by God in those already destined for him. The soul's assured salvation and the painful removal of what remains disordered belong together.

St.~Paul speaks of work tested by fire, with some work enduring and other work consumed, while the person is saved through fire (1~Cor 3:10--15). The passage first concerns building upon the foundation which is Jesus Christ, and the Church receives it among the scriptural witnesses to purification (CCC 1031). Its image teaches the gravity of what we build. Much that appears substantial can prove unable to endure the Lord's scrutiny. Fidelity must be founded upon Christ and formed by charity, rather than upon the impressive appearance which our work presents to others.

Benedict XVI considers the theological proposal that the purifying fire is the encounter with Christ himself, Judge and Savior (\emph{Spe salvi} 47). In this account his love reveals and burns away the falsehood which cannot endure communion with him. The soul desires the Savior whose love also exposes what needs purification. This explanation invites contemplation of an encounter in which judgment and healing come from Christ.

We can cooperate with purification now. Patient acceptance of unavoidable trials, prayer, works of mercy, and fitting penance can deepen charity and detach us from sin (CCC 1473). This does not justify seeking harm or treating another person's suffering casually. A difficult duty may expose impatience or resentment that was less apparent when life was easier. We can ask Christ for help, accept necessary assistance, and try to fulfill the duty with greater charity. The burden can then become an occasion to recognize and resist a fault.

The small renunciation can be especially revealing. We surrender the final word in an argument, the pleasure of repeating a humiliating story, the comfort purchased at another person's expense, or the demand to be thanked. These acts can seem slight, but they touch habits by which the self keeps returning to its own advantage. Their worth lies in charity. We undertake them because we want to love God more wholly and our neighbor more truthfully, accepting the loss that this purification requires.

The communion of saints extends our love to the faithful departed. The Church offers prayers in suffrage for them, above all the Eucharistic sacrifice, and commends works of charity and penance on their behalf (CCC 1032). These acts truly assist those being purified. They do more than preserve our memory of the dead. They belong to the communion in Christ through which his members can help one another even beyond death. Their destiny is already settled in grace; our suffrages serve their purification within that assured salvation.

Benedict relates this practice to the real interdependence of persons. Our lives enter one another's lives, and love in Christ does not become external to another person when death occurs (\emph{Spe salvi} 48). We should let that truth shape remembrance. We can have Mass offered for those we have loved, pray for people whose names are disappearing from earthly memory, and persevere after the first intensity of grief has passed. We can also pray for someone toward whom we still bear resentment, asking God to bring his saving work to fulfillment.

All Souls brings these duties before the Church each year. We entrust the dead to divine mercy and pray for their purification, leaving their judgment to Christ. A family remembrance can include having Mass offered and continuing prayer after the anniversary has passed. Such suffrage answers a need which earthly praise alone cannot meet. It also gives us reason to ask those close to us to pray for us when we die.

